\section{Background}

\subsection{Closure Cost Literature}

The economic impact of highway incidents has been studied primarily in the context of urban
traffic management, where incident clearance time on a congested freeway creates a queue
that imposes delay costs on following vehicles. \citet{Chin2004} developed an incident
cost estimation framework for FHWA that computed the secondary delay cost of incidents
on urban freeways as a function of incident severity, clearance time, and ambient V/C.
\citet{Taylor2006} extended this framework to rural freeways, finding that rural incidents
impose higher per-vehicle delay costs on average because traffic is more thinly spread
across fewer alternates. Neither framework addresses the case of complete corridor closure
with multi-hour duration and mandatory rerouting --- the dominant cost mechanism for
weather-driven mountain pass closures.

FHWA's incident cost estimation guide \citep{FHWA_incident2022} provides a national
incident database covering reported closures of 2+ hours on interstate routes. This
database is the primary data source for closure frequency in the current analysis. A
key limitation: the FHWA database undercounts rural incidents relative to urban incidents
by an estimated 30--40\%, because rural state DOTs have fewer automated incident detection
systems and lower reporting compliance for short-duration events. Our Donner cost estimate
is therefore conservative: actual closure frequency likely exceeds the 50 events/year
Caltrans figure when minor closures (30--90 minutes) are included.

\subsection{ATRI Cost Model}

The American Transportation Research Institute's operational cost model \citep{ATRI_costs2024}
provides the truck cost per hour used in this analysis. The 2024 figure is \$225 per
truck-hour for a fully loaded Class 8 truck, comprising:

\begin{itemize}
\item Driver wages and benefits: approximately \$0.67 per mile, annualized to \$73/hr at
55 mph average speed
\item Fuel: approximately \$0.52 per mile (diesel at \$3.90/gal, 7.5 mpg) = \$28.60/hr
\item Truck depreciation and finance: approximately \$0.18/mi = \$9.90/hr
\item Insurance: approximately \$0.12/mi = \$6.60/hr
\item Maintenance: approximately \$0.14/mi = \$7.70/hr
\item Schedule disruption premium: approximately \$100/hr for loads subject to just-in-time
delivery constraints or shipper detention fees
\end{itemize}

The \$225/hr figure is appropriate for T1 freight corridors where JIT scheduling is common.
For lower-tier corridors with bulk commodity freight (grain, coal, lumber), the relevant
figure is closer to \$150/hr, which omits the schedule disruption premium. We apply
\$225/hr uniformly in this analysis, with the conservative note that bulk-commodity
dominant corridors may have lower effective per-hour costs.

\subsection{Rerouting Penalty}

When a corridor closes, freight trucks must choose between waiting at the closure point
or rerouting via an alternate. The rerouting decision depends on expected closure duration.
For short closures ($<$2 hours), waiting is typically preferred. For long closures ($>$4
hours on a mountain pass), rerouting via a multi-hundred-mile alternate is preferred.
The rerouting cost is:

\begin{equation}
C_\text{reroute} = \$225/\text{hr} \times \frac{\text{detour miles}}{\text{avg speed
(mph)}} + C_\text{detention}
\end{equation}

where $C_\text{detention}$ is any appointment-related delay cost incurred at the
destination. For Donner (550-mile I-40 detour, average speed 60 mph including fuel
stops): rerouting cost = \$225 $\times$ 9.17 hr = \$2,063 per truck + detention.

\subsection{Waiting Penalty and JIT Supply Chains}

The alternative to rerouting --- waiting at the closure --- also imposes cost. At \$225/hr,
a truck waiting 18 hours at a Donner closure incurs \$4,050 in waiting cost plus driver
rest requirements. For JIT loads (perishables, automotive parts, retail distribution),
the delay may trigger shipper detention fees, missed unloading windows, or in extreme
cases (fresh produce), cargo spoilage.

\citet{Chin2004} estimated a ``waiting multiplier'' of 1.4--1.8 for JIT-sensitive loads:
the true economic cost of a waiting truck exceeds the ATRI truck-hour figure by 40--80\%
once shipper-side costs are included. We apply a waiting multiplier of 1.3 uniformly (conservative) in the closure cost model.

\subsection{V/C Ratio vs. Closure Cost: Low Correlation}

A preliminary analysis of V/C and annual closure cost across the ROUTE corpus reveals
a Pearson correlation coefficient of $r = 0.31$ (95\% CI: 0.18--0.43). This low correlation
confirms the intuition that the two metrics measure different phenomena. V/C measures
throughput stress on an operating corridor. Closure cost measures disruption cost when
the corridor ceases to operate. The most costly closures occur on corridors that are
\textit{not} congested under normal operations --- mountain passes, rural gulf coast
routes --- because these corridors have high isolation (B1) and high freight volume
relative to capacity, but low peak V/C because they operate below capacity except in
the rare events that close them entirely.
